Here we will talk about some efforts with sampler optimization. We will talk about the original sampler algorithm, the motivation behind our proposed rolling diffusion; some cost analysis for the naive variant of the algorithm; and how an adaptive variant can easily be derived by modifying the selection rule within the original sampler. 

\subsection{Original Sampling Algorithm}
The original DiffusionGemma sampler is based on the entropy-bound (EB) sampler first introduced in \citet{ben2026accelerated}, where the core idea is to commit tokens during the inference process based on the model's entropy. In short, given some entropy budget per step, the inference algorithm sorts all positions within a diffusion block according to entropy. 
Positions are then selected from lowest to highest until the entropy bound is satisfied. 
The selected positions are sampled. 

The base DiffusionGemma sampler combines this with a linear temperature schedule over the logits, as well as a noise injection: this can be considered almost a direct translation of the EB sampler into the uniform case. We include the algorithm below in Algorithm~\ref{alg:eb-base}.

\begin{algorithm}
\caption{EB Block Denoising}
\caption{EB Block Denoising}
\label{alg:eb-base}
\begin{algorithmic}[1] % The [1] enables line numbers. Remove it for no numbers.
    \REQUIRE Temperature schedule $\{\tau_T, \tau_{T-1}, \dots, \tau_{0}\}$, block size $d$, denoiser $D_{\theta}: |\mathcal{V}|^d \times \Delta_{|V|} \to \Delta_{|V|}$, stop condition: $f_{stop}: 
V^d \times \Delta_{d} \to \{ 0, 1 \}$, entropy budget $\gamma$
    \STATE $X_{cur} \sim U(V^d)$, $\mu_{cur} = \vec{0}_{|V| \times d}$
    \STATE $\tau_{cur} \gets \tau_{T}$
    \WHILE{not $f_{stop}(X_{cur}, \mu_{cur})$}
        \STATE $ \mu_{cur}  = \text{softmax} (D_{\theta}(X_{cur}, \mu_{cur}) / \tau_{cur})$
        \STATE sorted-indices, entropies = sort$(H(\text{Cat}(\mu_{new})))$
        \STATE update-tokens, update-pos $\gets \{\}, \{\}$
        \WHILE{running entropy $\gamma_{running} < \gamma$}
        \STATE $i \gets \text{sorted-indices.pop()}$ 
        \STATE $x \sim \text{cat}(\mu_{cur}[i])$
        \STATE $\gamma_{running} += \text{entropies}[i]$
        \ENDWHILE
    \STATE $X_{cur}[\text{update-pos}] \gets \text{update-tokens}$
    \STATE $X_{cur}[\neg \text{update-pos}] \gets U(V)$
    \ENDWHILE    
    \RETURN $X_{cur}$
\end{algorithmic}
\end{algorithm}

Given this algorithm for denoising each block, the full sampling algorithm is given in Algorithm~\ref{alg:block-ar}.

\begin{algorithm}
\caption{DiffusionGemma Block Autoregressive}
\label{alg:block-ar}
\begin{algorithmic}[1] % The [1] enables line numbers. Remove it for no numbers.
    \REQUIRE Prompt $X_{prompt}$, maximum sequence length $l_{max}$, block size $d$, denoiser $D_{\theta}$, encoder $E_{\theta}$, stop condition $f_{stop}$
    \STATE KV $\gets \text{Prefill}(X, l_{\max}, E_{\theta})$, $X_{cur} \gets X_{prompt}$
    \WHILE{not $f_{stop}(X_{cur})$}
        \STATE $X_{block} \gets $ EBBlockDenoise($D_{\theta}, X_{cur}, KV)$
        \STATE $X_{cur} \gets \text{concat}(X_{cur}, X_{block})$
        \STATE $\text{KV} \gets E_{\theta}(X_{cur}, KV)$
    \ENDWHILE
    \RETURN $X_{cur}$
\end{algorithmic}
\end{algorithm}
\subsection{Motivation: EB has an AR Bias}
While one defining characteristic of diffusion models is the fact that they can generate in any arbitrary order, it has been noted that adaptive sampling algorithms tend to "commit" to tokens in an autoregressive order \citep{ni2026flexibility}. To empirically examine this, we trace through the trajectories of the default EB sampler and record the order in which sub-blocks of tokens are committed in terms of percentiles -- i.e, by what fraction of the diffusion trajectory is is $j$ percentage of the $k$th block of tokens within a canvas committed to the final token, on average. 

\begin{figure}[h]
\begin{center}
\includegraphics[width=0.75\linewidth]{images/commitment_eb_gpqa.png}
\includegraphics[width=0.75\linewidth]{images/commitment_eb_lcb.png}
\includegraphics[width=0.75\linewidth]{images/commitment_eb_mmmu.png}
\end{center}
\caption{Commitment Frontiers for varying percentages of completions for sub-blocks of size $5$. For each task, we observe a consistent auto-regressive bias in that sub-blocks closer to the beginning of the canvas reach each visualized commitment rate earlier than the sub-blocks later in the canvas.}
\label{fig:feasibility}
\end{figure}

This observation leads to a source of inefficiency: even though the denoiser has already committed to the early tokens, it still spends subsequent forward passes on these tokens. It would be preferable to allocate these denoising steps to undetermined tokens later on in the sequence. Thus a natural optimization is to "roll forward" the denoising window by the number of tokens that are already stagnant / determined as opposed to waiting until the entire canvas is finished. 

\subsection{Rolling Diffusion: Simplest Form}
The core idea is straight forward: since tokens earlier on in the sequence are committed earlier, we should not spend forward passes on these. Since the tokens beyond the canvas are completely unrefined, we should instead reallocate denoising steps there. However, this introduces a few design choices / questions: 
\begin{enumerate}
    \item How do we decide when to move the canvas forward? 
    \item How do we handle the encoding of the tokens that are committed?
    \item How do we handle the temperature schedule, since tokens can now have different start times within the denoising window?
\end{enumerate}

We first will cover the simplest form, that just aims to demonstrate that it is possible to get the rolling approach to work.

\paragraph{Base Rolling Diffusion}
First, we observe that if we change the sorting rule in the EB algorithm in \ref{alg:eb-base} to be position based rather than entropy based, we arrive at a sampler that that accepts tokens strictly left to right. A reasonable rule for when to move the canvas forward is to define a "minimum token acceptance" threshold -- let us call this $T_{min}$. We can define a rule such that when the number of tokens accepted in the EB step exceeds $T_{min}$, we move the canvas forward by $T_{min}$. 

In this version of the algorithm, we aim to keep the core mechanics as close as possible to the original EB sampler as possible. First, we simply fix the tokens to be encoded by the autoregressive forward pass every single shift forward -- otherwise, this would result in the canvas attending to "dirty" KV cache that is contaminated by noisy tokens ahead, which cause the model to drift into OOD behavior. Second, to handle the temperature schedule, we simply interpret the EB schedule as applying to each position individually. A positions schedule begins when it first enters the canvas, and ends when it leaves the canvas by being committed.

\paragraph{Speculative Decoding with Rolling Diffusion}
A natural question that emerges is whether we can get more out of the autoregressive forward pass beyond just producing a clean KV cache from the tokens leaving the denoising window. One avenue is speculative decoding: 

